\documentclass[]{article}
% Added package start here
\usepackage{array}
\usepackage{geometry} % for margin of paper
\geometry{margin=1.34in}
\usepackage{fancyhdr} % for headers and footers
\usepackage{graphicx} % for image
\graphicspath{ {media/} }
\pagenumbering{gobble} % Delete page number
% End here
\usepackage{lmodern}
\usepackage{amssymb,amsmath}
\usepackage{ifxetex,ifluatex}
\usepackage{fixltx2e} % provides \textsubscript
\ifnum 0\ifxetex 1\fi\ifluatex 1\fi=0 % if pdftex
  \usepackage[T1]{fontenc}
  \usepackage[utf8]{inputenc}
\else % if luatex or xelatex
  \ifxetex
    \usepackage{mathspec}
  \else
    \usepackage{fontspec}
  \fi
  \defaultfontfeatures{Ligatures=TeX,Scale=MatchLowercase}
\fi
% use upquote if available, for straight quotes in verbatim environments
\IfFileExists{upquote.sty}{\usepackage{upquote}}{}
% use microtype if available
\IfFileExists{microtype.sty}{%
\usepackage{microtype}
\UseMicrotypeSet[protrusion]{basicmath} % disable protrusion for tt fonts
}{}
\usepackage{hyperref}
\hypersetup{unicode=true,
            pdfborder={0 0 0},
            breaklinks=true}
\urlstyle{same}  % don't use monospace font for urls
\IfFileExists{parskip.sty}{%
\usepackage{parskip}
}{% else
\setlength{\parindent}{0pt}
\setlength{\parskip}{6pt plus 2pt minus 1pt}
}
\setlength{\emergencystretch}{3em}  % prevent overfull lines
\providecommand{\tightlist}{%
  \setlength{\itemsep}{0pt}\setlength{\parskip}{0pt}}
\setcounter{secnumdepth}{0}
% Redefines (sub)paragraphs to behave more like sections
\ifx\paragraph\undefined\else
\let\oldparagraph\paragraph
\renewcommand{\paragraph}[1]{\oldparagraph{#1}\mbox{}}
\fi
\ifx\subparagraph\undefined\else
\let\oldsubparagraph\subparagraph
\renewcommand{\subparagraph}[1]{\oldsubparagraph{#1}\mbox{}}
\fi

\date{}

\begin{document}
% Added information start here
\pagestyle{fancy}
\renewcommand{\headrulewidth}{0pt}
% footer
  \fancyfoot[L]{\footnotesize This material is based upon work supported by the U.S. Department of Energy Office of Science, Advanced Scientific Computing Research and Biological and Environmental Research programs. \begin{flushright} \footnotesize Version 0.2, April 25, 2016 \end{flushright}}
  
\fancypagestyle{empty}{
% header
\fancyhead[C]{\Large {\textbf{What is Performance Portability}}\\ \Large {\textbf{for CSE Appilcations}}\\ \normalsize {The IDEAS Scientific Software Productivity Project} \\ \small {\href{https://ideas-productivity.org/resources/howtos/}{\emph{ideas-productivity.org/resources/howtos/}}}} 
  \fancyhead[R]{\includegraphics[width = 3 cm, height = 1.5 cm]{ideas_Whatis}}	
}
\thispagestyle{empty}
\textbf{\newline}
\textbf{\newline}
\textbf{\newline}

% Added information ends here.

\textbf{Portability:} An application code is portable if it can run on a
diverse set of platforms without needing significant modifications to
the source and can produce predictably similar output.

\textbf{Performance portability}: An application has portable
performance if in addition to running on diverse platforms it exhibits
similar accuracy, stability, and reliability across these platforms for
a given configuration. Moreover, the time to solution should reflect
efficient utilization of available computational resources on each
platform.

\textbf{Performance components in CSE applications}: The performance
of computational science and engineering (CSE) applications depends on
the model, discretization, numerical algorithm, input data, and
implementation. In roughly homogeneous platforms with similar basic
architecture, implementation details such as data structures,
synchronization, and other communication patterns dominate performance
portability of an application. With heterogeneity and diversity in
platform architecture, however, all the factors (the application model,
its discretization, numerical algorithm, input set, and implementation)
can impact performance portability.

\textbf{Performance factors}: The performance of a code can be
measured from many different perspectives. From the application science
perspective these include (1) quality of solution, (2) stability of the
code, and (3) time to solution. From the system perspective these
include (1) quality of machine utilization; (2) percentage of peak on
the machine, measured in operations per second (e.g., FLOPS/s),
bandwidth, or some other normalized value; and (3) throughput of the
machine.

\textbf{Constraints on performance:} Many CSE applications require
flexibility and composability because they address different physical
regimes either within the same simulation or in different instances of
simulations. Composability comes at the cost of raw performance and
plays both positive and negative roles in performance portability. On
the positive side, composability also allows flexibility and separation
of concerns, whereby different experts can focus on performance of
different aspects of the code. Also, composability encourages individual
components to be limited in scope, thus making it easier for the
compiler to analyze and optimize. On the negative side, increasing
composability decreases the compiler's ability to localize data use
between functions.

\textbf{Performance portability in the past:} \emph{Golden Age of
Vectorization} -- 1970-1980s -- with vector machines (mostly Cray).
\emph{Golden Age of MPI} -- 1995-2005 on systems with high-performance
networks. The dominant abstract machine model that enabled applications
to achieve portable performance was that of a powerful CPU with 2-3
levels of cache hierarchy and a generous amount of DRAM. The dominant
parallelization mode was the distributed-memory model. The finer details
differed, and maximizing performance on any one platform needed specific
optimization; but reasonable performance could be obtained across
platforms by programming to the general abstract machine model.

\textbf{Current impediments to performance portability:} At this time
the most common solution for performance portability is to maintain a
different version of the performance-critical sections of the
application code, for every machine variant. In different applications
this may range from a few kernels to most of the application. Some
applications may even require completely different organization of the
code and data structures. Known machine variants at present are multi-
or many-cores with or without accelerators, and whether the data
transfer occurs across PCI interface. An additional factor is the degree
of available SIMD or vectorization. Most often used programming
languages (C, C++, and Fortran) and parallel programming environments
(MPI and \href{http://openmp.org/wp/}{\emph{OpenMP}}) in CSE cannot
provide a solution at present, although MPI with the latest OpenMP
standards may provide a reasonable solution in future. One of the major
obstacles to performance portability is the diverse and conflicting set
of constraints on memory access patterns across devices. Contemporary
portable programming models address many-core parallelism (e.g.,
\href{http://openmp.org/wp/}{\emph{OpenMP}},
\href{https://www.khronos.org/opencl/}{\emph{OpenCL}},
\href{http://www.openacc.org}{\emph{OpenACC}}) but fail to address
memory access patterns.

No general solutions currently exist for achieving good performance
across all the variants. Individual application codes have achieved some
degree of performance portability through development of domain-specific
languages and custom back-ends for the target platforms, but they are at
best boutique solutions.

\textbf{Future possibilities:} Researchers at DOE laboratories and in
industry are exploring software abstractions based on standard C++
features to solve HPC performance and portability problems for current
and future platforms. The
\href{http://www.sciencedirect.com/science/article/pii/S0743731514001257}{\emph{Kokkos}}
and
\href{https://e-reports-ext.llnl.gov/pdf/782261.pdf}{\emph{RAJA{[}}}5{]}
C++ libraries, developed by Sandia National Laboratories and Lawrence
Livermore National Laboratory, respectively, attempt to enable
applications and domain libraries to achieve performance portability on
diverse many-core architectures by unifying abstractions for both
fine-grained data parallelism and memory access patterns. Nvidia is
taking a similar approach in the
\href{https://github.com/harrism/hemi}{\emph{HEMI}} library. Another
approach for achieving performance portability is using domain-specific
languages (DSLs) that map code to the different back-end languages such
as CUDA and OpenMP, for example
\href{http://ieeexplore.ieee.org/xpl/login.jsp?tp=\&arnumber=6495916\&url=http\%3A\%2F\%2Fieeexplore.ieee.org\%2Fxpls\%2Fabs_all.jsp\%3Farnumber\%3D6495916}{\emph{PyOP2}}.
These DSLs create a portable code across multicore architectures, but
they are limited to one specific problem domain.

\textbf{FAQ:}

\textbf{Q:} I use my code only on a local cluster, so why should I worry
about performance portability?

\textbf{A:} In the past, nodes of the cluster were relatively uniform,
and the parallel programming model was distributed programming. Now,
heterogeneity is coming even at the node level, so performance
portability is everyone's problem.

This document was prepared by Anshu Dubey with key contributions from
Todd Gamblin, Michael A. Heroux, Irina Demeshko, and Barry Smith.

\end{document}
